\section{Cut elimination in intuitionistic sequent calculus}
\label{sec:p07-intuitionistic-cut-elimination}

The intuitionistic cut-elimination theorem concerns the exact calculus
\ref{sheet:lj}.  Its sequents have a finite-multiset antecedent and exactly
one succedent formula.  The proof below removes both Cut and Mix, including
cuts whose formula is contracted, quantified, or introduced by implication
left.  It supplies the intuitionistic half of
\ref{thm:p07-cut-elimination}.  Gentzen's applications of the Hauptsatz
include consistency and the intuitionistic disjunction property
\cite[IV, \S1.1 and \S1.3]{Gentzen1935b}.  The construction here is stated for this
handbook's particular structural rules and literal substitution convention;
the exactly-one-succedent restriction is part of that specified calculus.

\subsection{The exact calculus and the induction objects}
\label{sec:p07-lj-cut-objects}

\begin{hbconvention}{def:p07-lj-cut-scope}{Logical and structural scope}
Throughout this chapter, a derivation means a finite tree using the rules
of \ref{sheet:lj}.  A \emph{cut-free} derivation uses neither
\ref{rule:lj-cut} nor \ref{rule:lj-mix}.  Its formula language has
\(\bot,\land,\lor,\to,\forall,\exists\), no equality, and defined
negation \(\neg A\defeq A\to\bot\).  Terms and substitution are those of
\ref{def:p02-fol-capture-avoiding-substitution}.  The initial identity is
the uncontexted sequent \(A\Rightarrow A\); additional antecedent formulas
require left weakening.  Exchange is already accounted for by multisets.
There is no right weakening, right contraction, or empty-succedent sequent.
No nonlogical initial sequents are included.  The object-variable alphabet
is the countably infinite alphabet of \ref{def:p02-fol-signature}, so a
fresh variable can always be chosen outside a finite collection of variables.
\end{hbconvention}

\begin{hbdefinition}{def:p07-lj-alpha-proof}{Proofs on alpha-equivalence classes}
Write \([A]_\alpha\) for the class of a literal formula \(A\) under the
bound-variable renaming relation \(\equiv_\alpha\) of
\ref{def:p02-fol-alpha-equivalence}. Equality of these classes refers to
that same syntactic relation, not semantic equivalence. An \emph{alpha-proof} is a finite
proof tree whose formula occurrences are such classes and whose rule
instances are the images, modulo alpha-equivalence, of the rules in
\ref{sheet:lj}.  Antecedents remain multisets of occurrences and succedents
still contain exactly one occurrence.  Replacing a label by another
representative of its class is a choice of notation in this auxiliary
representation, not an additional rule of the literal calculus.

Define formula rank recursively by
\[
 \begin{aligned}
 \rho(P(\vec t))&=0,\qquad \rho(\bot)=1,\\
 \rho(B\circ C)&=1+\rho(B)+\rho(C),\\
 \rho(QxB)&=1+\rho(B).
 \end{aligned}
\]
where \(\circ\in\{\land,\lor,\to\}\) and
\(Q\in\{\forall,\exists\}\).  Rank does not count symbols inside terms.
Proof height \(h(D)\) is zero for an initial sequent and one plus the
maximum premise height at an inference.  Alpha-equivalent formulas have
the same rank.  These are the measure conventions of
\ref{entry:p07-source-15} restricted to the present language.

A finite simultaneous substitution \(\theta\) assigns terms to
finitely many variables and fixes every other variable.  It replaces
free occurrences simultaneously: a term inserted for a variable is
not itself subjected again to \(\theta\) during that operation.
Bound variables are renamed when needed to avoid capture.  The
resulting alpha-class is independent of those fresh choices, as proved
in the next lemma.  Write \(B\theta\) and \(t\theta\) for its action
on a formula and on a term, respectively.
\end{hbdefinition}

\begin{hblemma}{thm:p07-lj-alpha-substitution}{proved=annex:p7.lj.alpha-substitution}{Substitution and fresh eigenparameters}
Every literal derivation projects to an alpha-proof of the same height.
For a cut-free alpha-proof, simultaneous capture-avoiding substitution of
finitely many terms for free variables preserves derivability and does not
increase height.  The eigenparameters of its quantifier inferences can
first be chosen outside any prescribed finite set, consistently with their
premise subtrees.  Formula rank is unchanged by these substitutions.
\end{hblemma}

\begin{hbwarning}{def:p07-lj-literal-height}{Why the auxiliary representation is needed}
Height preservation in the preceding lemma is a statement about
alpha-proofs.  It is not a claim about independently changing literal
representatives in the original sheet.  For example,
\(\forall xP(x)\Rightarrow\forall xP(x)\) is a height-zero identity,
whereas \(\forall xP(x)\Rightarrow\forall zP(z)\), with different bound
variable names, is not literally an identity instance.  Also, the fixed
least-fresh convention for capture avoidance can make two successive
substitutions agree only up to alpha-equivalence.  The final literal
realization theorem below addresses both points without adding a rule.
\end{hbwarning}

\begin{hbdefinition}{def:p07-lj-marked-multicut}{Selected occurrences and multicut}
Let \(D\) prove \(\Gamma\Rightarrow A\), and let \(E\) prove a sequent
whose antecedent contains a selected list of \(n\) occurrences of \(A\).
Write it as \(A^n,\Pi\Rightarrow C\), with the selected occurrences in
\(A^n\) and every other occurrence in \(\Pi\).  The multiset \(\Pi\)
may itself contain occurrences of the same formula \(A\); these remain
unselected.  The desired result is \(\Gamma,\Pi\Rightarrow C\), with
one copy of the whole multiset \(\Gamma\).

The selection is bookkeeping on occurrences.  An occurrence introduced by
a rule is principal for this operation only if it is selected.  A newly
exposed constituent or discharged assumption is unselected even if it
happens to be the same formula as \(A\).  In the induction we also allow
\(n=0\): its conclusion is obtained simply by weakening \(\Gamma\)
into \(E\).  The sheet's actual Mix rule still requires \(n\geq1\).
\end{hbdefinition}

\begin{hblemma}{thm:p07-lj-context-collapse}{proved=annex:p7.lj.context-collapse}{Collapse of repeated supplied contexts}
If a cut-free proof has endsequent \(\Gamma^m,\Pi\Rightarrow C\), with
\(m\geq1\) designated copies of the finite multiset \(\Gamma\), left
contraction gives a cut-free proof of \(\Gamma,\Pi\Rightarrow C\).
Every occurrence assigned to \(\Pi\) is retained, including occurrences
whose formula is also in \(\Gamma\).  The assertion holds for literal
proofs and for alpha-proofs.
\end{hblemma}

\begin{hblemma}{thm:p07-lj-bottom-replacement}{proved=annex:p7.lj.bottom-replacement}{Cut-free replacement of a bottom succedent}
If \(\Gamma\Rightarrow\bot\) has a cut-free alpha-proof, then
\(\Gamma\Rightarrow C\) has a cut-free alpha-proof for every formula
\(C\).  This construction uses bottom left and left inferences; it does
not use a right weakening rule.
\end{hblemma}

\subsection{The strengthened reduction theorem}
\label{sec:p07-lj-marked-reduction}

\begin{hbtheorem}{thm:p07-lj-marked-multicut}{proved=annex:p7.lj.marked-multicut}{Marked multicut admissibility}
For cut-free alpha-proofs \(D\) of \(\Gamma\Rightarrow A\) and \(E\)
of \(A^n,\Pi\Rightarrow C\), with the occurrence convention of
\ref{def:p07-lj-marked-multicut}, there is a cut-free alpha-proof of
\(\Gamma,\Pi\Rightarrow C\).  The construction terminates by
lexicographic induction on
\[
                \bigl(\rho(A),h(D)+h(E)\bigr).
\]
The zero-selection case is a direct weakening construction and makes no
recursive call.  Contraction may increase the number of selected
occurrences in a premise; it still decreases the second induction
coordinate.  No bound on the height of the resulting proof is asserted.
\end{hbtheorem}

\begin{hbmethod}{def:p07-lj-reduction-procedure}{How to apply the reduction}
Start with cut-free premises and carry the selected occurrences backwards
through the final inference of each proof.  Initial sequents finish the
construction directly.  If the last inference of the left proof does not
introduce its succedent, move the operation to its relevant premises.  If
the left succedent is introduced but the last right inference does not
introduce a selected antecedent occurrence, move the operation to the right
premises.  Restore the same inference afterwards and contract repeated
supplied contexts where its premises split.

When both final inferences introduce the selected formula, first eliminate
every residual selected occurrence from the relevant right premises,
using the original left proof and the strictly shorter right premise
proofs.  Only then perform the cuts on proper constituents, or on an
instance of a quantified matrix.  These latter calls have smaller rank,
regardless of their proof heights.  For a contracted selected occurrence,
select both predecessor occurrences.  For an eigenvariable inference,
freshen its parameter before moving a proof with additional free variables
across it.  The annex gives every structural, connective and quantifier
case of this procedure.
\end{hbmethod}

\subsection{Returning to literal formulas}
\label{sec:p07-lj-literal-realization}

\begin{hblemma}{thm:p07-lj-alpha-identity}{proved=annex:p7.lj.alpha-identity}{Literal identity between alpha-equivalent formulas}
If literal formulas \(A\) and \(A'\) are alpha-equivalent, there is a
cut-free literal derivation of \(A\Rightarrow A'\) in \ref{sheet:lj}.
The statement holds in both directions and is proved by formula-rank
induction.  It does not assert height zero.
\end{hblemma}

\begin{hblemma}{thm:p07-lj-literal-realization}{proved=annex:p7.lj.literal-realization}{Literal realization of a cut-free alpha-proof}
Given a cut-free alpha-proof, choose a literal representative for each
formula occurrence of its endsequent.  The exact resulting literal
sequent has a cut-free derivation in \ref{sheet:lj}.  In shared contexts
the chosen representatives are used consistently in the premise
derivations.  No height-preservation claim is required or made.
\end{hblemma}

\begin{hbtheorem}{thm:p07-lj-exact-cut-elimination}{proved=annex:p7.lj.exact-cut-elimination}{Cut and Mix elimination for the exact intuitionistic sheet}
Every derivation in \ref{sheet:lj} can be transformed into a derivation
using neither Cut nor Mix and having the same literal endsequent.
The theorem includes first-order formulas, finite-multiset antecedents,
left contraction, the sheet's split-context implication-left rule,
uncontexted identity, bottom left, and the stated eigenvariable and
capture-avoiding substitution conventions.
\end{hbtheorem}

\subsection{Analyticity, consistency and witnesses}
\label{sec:p07-lj-cut-consequences}

\begin{hbdefinition}{def:p07-lj-instance-subformula}{Subformula instances}
For a finite collection of formulas, its \emph{subformula-instance
closure} is the smallest collection of alpha-classes containing their
classes and closed under taking immediate connective constituents,
taking the matrix of any quantified representative, and making finite capture-avoiding
substitutions of terms for free variables.  This closure permits
instances of a matrix whose bound variable has been opened by a
quantifier rule.  It places no finite bound on the available terms.
For a purely propositional language, use the ordinary literal subformula
closure instead.
\end{hbdefinition}

\begin{hbcorollary}{thm:p07-lj-analyticity}{proved=annex:p7.lj.analyticity}{Exact cut-free subformula conclusions}
Every formula occurrence in a cut-free first-order derivation belongs,
modulo alpha-equivalence, to the subformula-instance closure of its
endsequent.  In a propositional derivation every formula is an ordinary
literal subformula of an endsequent formula.  Consequently every
derivable sequent has a proof with the corresponding property.
\end{hbcorollary}

\begin{hbcorollary}{thm:p07-lj-consistency}{proved=annex:p7.lj.consistency}{Underivability of bottom from no assumptions}
There is no derivation of \(\Rightarrow\bot\) in \ref{sheet:lj}.
This is the consistency conclusion for its exactly-one-succedent syntax.
\end{hbcorollary}

\begin{hbcorollary}{thm:p07-lj-disjunction-property}{proved=annex:p7.lj.disjunction-property}{Disjunction property for theorems}
If \(\Rightarrow B\lor C\) is derivable in \ref{sheet:lj}, then at
least one of \(\Rightarrow B\) and \(\Rightarrow C\) is derivable.
Each can be taken cut-free when it exists.  The empty-antecedent
hypothesis is part of the assertion.
\end{hbcorollary}

\begin{hbcorollary}{thm:p07-lj-existence-property}{proved=annex:p7.lj.existence-property}{Existence property and the ground-term condition}
If \(\Rightarrow\exists xB\) is derivable in \ref{sheet:lj}, there is
a term \(t\) and a cut-free derivation of
\(\Rightarrow B[t/x]\).  The term may contain free variables.
If \(\exists xB\) is closed and the signature contains a ground term,
one can choose a ground witness term.  Without a ground term in the
signature, the latter conclusion is not available in that signature.
\end{hbcorollary}

\begin{hbexample}{def:p07-lj-parameter-witness}{A witness in a language with no closed terms}
Take a signature with a unary predicate \(P\) and no constants or
function symbols.  Identity gives \(P(y)\Rightarrow P(y)\);
implication right gives \(\Rightarrow P(y)\to P(y)\); existential
right with witness \(y\) gives
\(\Rightarrow\exists x(P(x)\to P(x))\).  This is already a cut-free
derivation, but the signature has no ground terms.  Adding a fresh
constant would change the signature and is an explicit separate option.
\end{hbexample}

\begin{hbcorollary}{thm:p07-lj-quantifier-free-conservativity}{proved=annex:p7.lj.quantifier-free-conservativity}{Conservativity on quantifier-free sequents}
If a quantifier-free sequent is derivable in the first-order sheet, it
has a derivation using only its quantifier-free logical and structural
rules.  In that derivation every formula is a literal subformula of an
endsequent formula.  Its first-order atoms may therefore be treated as
propositional atoms for that derivation.
\end{hbcorollary}

\begin{proofboundary}{def:p07-lj-cut-boundary}{Limits of the result}
The proof supplies no polynomial size bound for cut elimination and no
finite enumeration of all first-order substitution instances.  It does
not extend the theorem to equality rules or arbitrary nonlogical axiom
sequents.  Defined negation is covered by the implication and bottom
cases.  Classical multiple-succedent interpolation, midsequent and
Herbrand results in the surrounding chapters retain their own statements
and proofs.  The disjunction property above concerns theorems: from the
assumption \(P\lor Q\), identity proves that disjunction, while neither
\(P\lor Q\Rightarrow P\) nor \(P\lor Q\Rightarrow Q\) is derivable
for distinct propositional atoms.  The annex verifies this last failure
directly by cut-free last-inference analysis.
\end{proofboundary}
